\section{Overview of Stream Processing}

\subsection{Batch vs.\ Stream Processing}

Data processing systems can be broadly categorised by when computation occurs
relative to data arrival:

\begin{itemize}
  \item \textbf{Batch processing} — data is collected over a time window, stored, and then
        processed in a single large job. Systems like Apache Hadoop MapReduce and scheduled
        SQL pipelines fall into this category. Batch processing is well-suited to historical
        reports and large-scale transformations where latency is measured in hours or days.
  \item \textbf{Stream processing} — computation is performed \emph{continuously} as each record
        arrives. Results are produced with latency measured in milliseconds to seconds.
        This model is essential for real-time dashboards, fraud detection, IoT monitoring,
        and any use case where acting on stale data is unacceptable.
\end{itemize}

\subsection{Key Challenges in Stream Processing}

Building a correct, efficient stream-processing system is significantly harder than
batch processing, due to several inherent challenges:

\begin{itemize}
  \item \textbf{Unbounded data} — the stream has no defined end; the system must process records
        indefinitely without accumulating unbounded state.
  \item \textbf{Out-of-order events} — network delays, retries, and clock skew mean that events
        frequently arrive out of their logical order. A system must decide when it is safe to
        emit a windowed result even if some late events may still be in transit.
  \item \textbf{Stateful computation} — aggregations, joins, and deduplication require maintaining
        keyed state across events. This state must be fault-tolerant and scalable.
  \item \textbf{Fault tolerance} — any node in the cluster can fail at any time.
        The system must recover to a consistent state without losing records or producing duplicates.
  \item \textbf{Delivery semantics} — depending on the use case, the system must guarantee
        \emph{at-most-once} (may lose), \emph{at-least-once} (may duplicate), or
        \emph{exactly-once} (no loss, no duplicate) processing.
\end{itemize}

\subsection{Two Leading Engines}

This chapter examines the two most widely adopted open-source stream-processing frameworks:

\begin{description}
  \item[Apache Spark Structured Streaming] — extends Spark's batch DataFrame API to streaming data
        via a \emph{micro-batch} execution model, with an experimental \emph{continuous processing}
        (real-time) mode for lower latency.
  \item[Apache Flink] — a purpose-built, true-streaming engine that processes events one at a time
        using a \emph{dataflow graph} model, with first-class support for event time, watermarks,
        and exactly-once semantics via the Chandy-Lamport checkpointing algorithm.
\end{description}
